% 第4回　文献講読と批判的思考（記入例・模範解答）
\session{4}{1}{10月21日（水）}{文献講読と批判的思考}{AI要約と対話}{}

\goals{
  \item 論文の構造（問い、方法、結果、考察、限界）を押さえて読める
  \item AI要約の使いどころと限界を説明できる。要約は原文の代わりにならない
  \item 主張と根拠の対応、データの範囲、別の解釈の可能性を確かめる批判的な読み方ができる
}

\afillin{読む文献（著者・年・題名）}{架空論文Ｘ「都市部の若年層における自動車購入前の情報探索と販売店訪問」}

\begin{caution}{入力するときの注意}
論文の全文を入力する場合は、大学のAI利用規程と、その資料の利用条件を確認する。要旨や必要な部分だけを使えば十分なことが多い。
\end{caution}

\work[80mm]{ワーク1}{AI要約と原文を照らし合わせる}
\begin{promptbox}[構造をつかむ]
次の論文について、「問い」「方法」「主な結果」「著者が認めている限界」をそれぞれ2文以内でまとめてください。書かれていない項目は「記載なし」としてください。（論文の要旨や該当部分を貼る）
\end{promptbox}
\begin{wstabh}{colspec={Q[l,wd=17mm,m] X[l,m] X[l,m] X[l,m]},row{2-Z}={ht=17mm}}
  & AIの要約（要点） & 原文で確認した内容（頁） & 要約の誤り・抜け \\
  問い & \af{若者はなぜ販売店に行かなくなったのか} & \af{購入前のオンライン情報探索の量と、訪問する店舗数・回数の関係（p.2）} & \af{「なぜ行かなくなったか」は論文の問いではない。要約が問いを広げている} \\
  方法 & \af{若年層へのアンケート調査} & \af{都市部の20〜34歳を対象にしたWeb調査。過去3年に車を購入した人に限定（p.4）} & \af{「購入経験者に限定」が抜けている。結果の範囲に関わる重要な点} \\
  主な結果 & \af{情報を多く集める人ほど店に行かない} & \af{情報探索が多い人ほど訪問店舗数は少ないが、1回の訪問での商談時間は長い（p.8）} & \af{後半（商談が長い）が抜け、「店に行かない」と言い過ぎている} \\
  限界 & \af{記載なし} & \af{Web調査のため回答者に偏りがある。一時点の調査で因果は言えない（p.12）} & \af{原文には限界の節がある。AIが見落とした} \\
\end{wstabh}

\work[70mm]{ワーク2}{最も強い反論に、原文の根拠で答える}
\begin{promptbox}[反論を引き出す]
この論文の主張「（主張）」に対して、最も強い反論を3つ挙げてください。それぞれについて、どのようなデータがあれば検証できるかも示してください。
\end{promptbox}
\begin{wstabh}{colspec={X[3,l,m] Q[c,wd=15mm,m] X[3,l,m]},row{2-Z}={ht=16mm}}
  AIが挙げた反論 & 答えられる\newline ○△× & 原文の根拠（頁）・答えられない理由 \\
  \af{車を買った人だけを調べているので、買わなかった人の来店の変化はわからない} & \ans{×} & \af{原文も対象を購入経験者に限っている（p.4）。この反論には答えられず、論文の範囲の外} \\
  \af{情報を多く集める人は、もともと車に詳しい人ではないか（別の要因）} & \ans{△} & \af{車への関心度で分けた分析がある（p.9）が、関心の測り方が1問だけで弱い} \\
  \af{オンラインの情報源の種類（動画・口コミ・公式サイト）を区別していない} & \ans{○} & \af{情報源別の集計がある（表3、p.7）。動画の利用が多い人で訪問店舗数が少ない} \\
\end{wstabh}

\work{ワーク3}{未解決の要素}
\atwoboxes{論文が答えていないこと}{自分の問いとの隙間}{32mm}
{・車を買わなかった人の来店はどう変わったか\newline ・訪問店舗数が減った理由（情報で絞り込んだのか、行くのが負担なのか）\newline ・地方と都市の違い}
{論文は「都市部」全体。名古屋都市圏のように車の利用が多い地域では、結果が違う可能性がある。また、来店の「回数」だけでなく、来店の目的（見る・試乗・商談）の変化を見たい。}

\work{共有}{グループで出た意見・他の人の読み方から学んだこと}
\alines{3}{AIの要約は「言い過ぎ」になりやすいという点で、全員の意見が一致した。他の人は、表の注（対象・期間）を最初に読んで、結果の範囲を確かめていた。次からまねしたい。}

\begin{nexttime}
  \item 読んだ文献について「主張」「根拠」「限界」を各100字程度でまとめる
  \item AIとの議論で見えた未解決の要素を1つ書き、自分の問いとどうつながるかを考える
\end{nexttime}
\abox[主張（100字程度）]{24mm}{都市部の若年層では、購入前にオンラインで多くの情報を集める人ほど、訪問する販売店の数は少ない。ただし、来店1回あたりの商談時間は長く、来店は「比べる場」から「決める場」に変わりつつある。}
\abox[根拠（100字程度）]{24mm}{過去3年に車を購入した都市部の20〜34歳へのWeb調査で、情報探索の量と訪問店舗数・商談時間の関係を集計した。情報源別に見ると、動画の利用が多い人ほど訪問店舗数が少なかった。}
\abox[限界（100字程度）]{24mm}{対象が購入経験者に限られ、買わなかった人の行動はわからない。Web調査で回答者に偏りがあり、一時点の調査なので因果関係は言えない。車への関心度の測り方も1問だけで弱い。}
\abox[未解決の要素と、自分の問いとのつながり]{24mm}{来店の「目的」がどう変わったかは調べられていない。自分の問いでは、名古屋都市圏の20代を対象に、来店前に何を調べ、何のために店に行くのかを聞くことで、この隙間を埋められる。}
\aimemoA{論文の構造の要約と、主張への反論の生成}
{「問い・方法・主な結果・限界を2文以内で。書かれていない項目は記載なしと」「最も強い反論を3つ」}
{要約は読む前の地図としてだけ使い、表の内容はすべて原文から書き直した。反論3つのうち2つを採用。}
{要約の各項目を原文の該当ページと照合した（誤りと抜けを表に記録）}

\begin{point}
  \item 記入例で使った論文は、この冊子のための\textbf{架空の論文}。授業では、各自が第3回で入手した文献を使う。
  \item 評価の中心は「要約の誤り・抜け」の列。AI要約が問いを広げる、条件（対象の限定）を落とす、結果を言い過ぎる、限界を落とす、という典型的なずれを見つけられているかを見る。
  \item 原文の頁を書かせることで、原文を実際に読んだかを確かめられる。頁のない記入は、原文を読まずにAIの要約を写している可能性がある。
  \item 「答えられない反論」が出てくるのは悪いことではない。それが論文の範囲の外だと言えれば十分で、自分の問いの出発点になる。
\end{point}
